\chapter{Reproducibility}\label{app:repro}

% This appendix is dense with monospace paths, which TeX will not hyphenate. At
% the default tolerance it prefers to run a line past the right margin over
% breaking early, and five paragraphs here did exactly that. Raising the
% tolerance for the appendix trades a little inter-word stretch for text that
% stays inside the text block.
% Path-heavy tables and URLs need looser lines than the body, but \sloppy
% over a whole appendix buys that with visibly loose interword spacing on
% ordinary prose. \emergencystretch is only drawn on when a paragraph would
% otherwise overfull, so normal paragraphs set normally.
\setlength{\emergencystretch}{2em}

\section{Figure-to-code map}

Every computed figure in this thesis is built by a script in this folder's
\texttt{tools/}: \texttt{make\_thesis\_figures.py} and one
\texttt{fig\_chNN.py} per results chapter, all styled through
\texttt{tools/figstyle.py} (one palette and layout module, so every figure
looks like it came from the same hand; text is embedded as TrueType). The
closed-form panels are recomputed from the formulas at build time; the
measured panels read the frozen research outputs under
\texttt{research/outputs/frozen/} directly --- the same files
the generated tables read, so a figure and the table beside it cannot
disagree about a number. Randomness, where used, is seeded. Seven figures are
reproduced from the literature with attribution in their captions and are not
built here (\texttt{fig\_song\_dog}, \texttt{fig\_song\_sde},
\texttt{fig\_bachtis\_cascade}, \texttt{fig\_bonnaire\_memorize},
\texttt{fig\_biroli\_regimes}, \texttt{fig\_mezard\_bp\_messages},
\texttt{fig\_unet\_architecture}); two are Ti\emph{k}Z schematics drawn in
the source.

\begin{center}
\begin{minipage}{\linewidth}\small\centering
\captionof{table}[Figure-to-code map]{Every computed figure, the file it is
committed as, and the script that builds it. Scripts marked \emph{frozen}
read the frozen research outputs; the others recompute closed forms.}
\label{tab:figure-map}
\begin{tabular}{@{}lll@{}}
\toprule
Figure & File & Built by \\
\midrule
\ref{fig:sm-corruption} & \texttt{fig\_forward\_corruption.pdf} & \texttt{make\_thesis\_figures}, closed forms \\
\ref{fig:toymodel} & \texttt{fig\_marginal\_blindness.pdf} & \texttt{make\_thesis\_figures}, closed forms \\
\ref{fig:ring-model-dynamics} & \texttt{fig\_ring\_model\_dynamics.pdf} & \texttt{make\_thesis\_figures}, simulation \\
\ref{fig:ring-model-channel} & \texttt{fig\_ring\_model\_channel.pdf} & \texttt{make\_thesis\_figures}, simulation \\
\ref{fig:ring-gauge} & \texttt{fig\_ring\_gauge.pdf} & \texttt{make\_thesis\_figures}, simulation \\
\ref{fig:g-eigen-relax} & \texttt{fig\_eigen\_relax.pdf} & \texttt{fig\_ch05}, closed forms \\
\ref{fig:g-sparsity-lifecycle} & \texttt{fig\_sparsity\_lifecycle.pdf} & \texttt{fig\_ch05}, closed forms \\
\ref{fig:g-markov-survival} & \texttt{fig\_markov\_survival.pdf} & \texttt{fig\_ch05}, closed forms \\
\ref{fig:truncation} & \texttt{fig\_truncation.pdf} & \texttt{fig\_ch06}, closed forms \\
\ref{fig:l-nonlinear} & \texttt{fig\_laplace\_nonlinear.pdf} & \texttt{fig\_ch07}, quadrature \\
\ref{fig:l-metrics} & \texttt{fig\_laplace\_metrics.pdf} & \texttt{fig\_ch07}, exp.~02, frozen \\
\ref{fig:bp-closure} & \texttt{fig\_laplace\_closure.pdf} & \texttt{fig\_ch07}, exp.~02, frozen \\
\ref{fig:em-grid} & \texttt{fig\_em\_grid.pdf} & \texttt{fig\_ch08}, exp.~01/18, frozen \\
\ref{fig:em-diagnostics} & \texttt{fig\_em\_diagnostics.pdf} & \texttt{fig\_ch08}, exps.~18 and 06, frozen \\
\ref{fig:em-innovation} & \texttt{fig\_em\_innovation.pdf} & \texttt{fig\_ch08}, exp.~06, frozen \\
\ref{fig:em-two-rates} & \texttt{fig\_em\_two\_rates.pdf} & \texttt{fig\_ch09}, exp.~27, frozen \\
\ref{fig:em-value-of-structure} & \texttt{fig\_value\_of\_structure.pdf} & \texttt{fig\_ch10}, exp.~07/31, frozen \\
\ref{fig:em-screening} & \texttt{fig\_screening.pdf} & \texttt{fig\_ch10}, exp.~31 screen, frozen \\
\ref{fig:em-capacity} & \texttt{fig\_capacity.pdf} & \texttt{fig\_ch10}, exp.~32, frozen \\
\ref{fig:em-nonmarkov} & \texttt{fig\_nonmarkov.pdf} & \texttt{fig\_ch10}, exp.~21 clean rerun \\
\ref{fig:em-information} & \texttt{fig\_information.pdf} & \texttt{fig\_ch10}, exp.~22, frozen \\
\bottomrule
\end{tabular}
\end{minipage}
\end{center}

\section[The learning experiments]{The learning experiments of
Chapters~\ref{ch:em-method} and~\ref{ch:em-results}}\label{sec:repro-learning}

The closed-form and quadrature studies above are deterministic and run on a
laptop. The estimator comparisons are neither, and they are gated
differently.

\paragraph{Provenance.} Runs deployed by \texttt{hpc/deploy\_clean.sh} are
certified: the script refuses to build from a dirty tree, ships a
\texttt{git archive} of the commit rather than the working tree, records the
SHA-256 of that tarball, and extracts on the cluster under the commit id so that
a run cannot overwrite or edit its own source. The motivating failure is
preserved in that script's header: an earlier two-step stamp-then-rsync path
produced \texttt{outputs/exp\_12\_scaled/}, whose recorded commit cannot have
run the recorded command, and whose numbers are therefore not recoverable and
are not cited in this thesis. \texttt{tools/provenance\_gate.py} intersects a
run's recorded dirty-file list with the import closure of its own experiment
module, so that uncommitted work elsewhere in the tree is distinguished from
uncommitted work on the code path that produced the numbers. That distinction
does the work below, because the two are not equally serious.

\paragraph{How each run is certified.}\label{sec:repro-provenance}
Not every experiment here was run after
\texttt{deploy\_clean.sh} existed, so four routes appear, in decreasing
order of strength. Nothing cited in this thesis hides which route it took. The development record behind each
route, which deployments failed, and why each check was added, is in
\compendium{}.

\emph{Certified.} \texttt{exp\_31}, the architecture screen and the
confirmatory comparison, and \texttt{exp\_32}, capacity, were deployed by the
script and pass \texttt{require\_clean} unmodified. Every number in
Sections~\ref{sec:em-architecture} and~\ref{sec:em-capacity} comes from them.

\emph{Legacy digest, empty intersection.} \texttt{exp\_07}, the headline
efficiency comparison, its validation-selection audit and its budget
calibration, predates the script. Its records carry a commit, a host, a job
id, a resolved configuration and a dirty-file list, but no
archive digest, so the plain gate refuses it. What makes it citable is
narrower and checkable: the intersection of each run's recorded dirty list with
its own import closure is \emph{empty}, so no uncommitted edit could have
reached the code that produced the numbers.
\texttt{make\_tab\_efficiency.py} recomputes that intersection on every run
and refuses if it is ever non-empty, not resting on a check performed
once by hand. The one \texttt{exp\_07} row where the intersection was
\emph{not} empty, $\nseq = 8192$, whose dirty list contained the
experiment, the frozen configuration and \texttt{src/em.py}, is withdrawn
and appears nowhere in this thesis.

\emph{Recorded configuration only.} The validation and diagnostic runs ---
\texttt{exp\_01} and \texttt{exp\_18} (grid validation and diagnostics),
\texttt{exp\_02} (the Gaussian-closure sweep), \texttt{exp\_06} (parameter
recovery), \texttt{exp\_22} (the Fisher-information curves of
Section~\ref{sec:em-shape-information}) and \texttt{exp\_27} (the settling
times of Section~\ref{sec:em-stopping-times}) --- predate even the commit
recording: their \texttt{params} files carry the resolved configuration,
host, timestamp and, where run under SLURM, the job id, but an empty revision
field. They are reproducible by rerunning the committed experiment code under
the recorded configuration (the quadrature studies are deterministic; the
rest record their seeds), and the claims that rest on them are stated at
matching strength in the chapters --- measured curves and orderings, not
certified point values. This is the weakest of the four routes and is
disclosed as such.

\emph{Certified after rerun.} \texttt{exp\_21}, behind
Table~\ref{tab:nonmarkov} and Figure~\ref{fig:em-nonmarkov}, is reported from
\texttt{outputs/frozen/exp\_21\_clean}: fifteen cells from two git-archive
deployments, every one with an empty dirty list. It replaces an earlier run
whose recorded commit did not reconstruct the program, and the Gaussian arm
reproduces that run on every number quoted here, to the precision printed.
Because the Laplace and Gaussian arms were deployed separately, the rerun spans
two commits, so \texttt{make\_tab\_nonmarkov.py} refuses unless no single cell
pools commits \emph{and} every file in the experiment's sixteen-file import
closure is byte-identical between them. Both hold, and the check runs on every
regeneration.

\paragraph{Generated numbers.} No efficiency, structured-baseline or
aggregation number in this thesis is typed. Each is defined once as a macro in
\texttt{thesis/sections/}, written there by one of the thirteen generators
below, each reading the recorded outputs directly. Nine live under
\texttt{research/tools/}; three are thesis-side scripts that emit the macros
their own captions quote (\texttt{fig\_ch10},
\texttt{make\_laplace\_oracle} and \texttt{make\_build\_manifest}).
\texttt{make\_ref\_energy} is the fourteenth name and writes no macro at all:
it recovers the per-level reference energies of the frozen test bundle that
the pooled estimand needs, and the others read what it produces.

\begin{center}\small
\begin{tabular}{@{}ll@{}}
\texttt{make\_tab\_efficiency} & \texttt{make\_aggregation\_robustness} \\
\texttt{make\_tab\_structured} & \texttt{make\_aggregation\_structured} \\
\texttt{make\_tab\_screening}  & \texttt{make\_convergence\_numbers} \\
\texttt{make\_tab\_capacity}   & \texttt{make\_grid\_diagnostics} \\
\texttt{make\_tab\_nonmarkov}  & \texttt{make\_ref\_energy} \\
\texttt{fig\_ch10}             & \texttt{make\_laplace\_oracle} \\
\texttt{make\_build\_manifest} & \\
\end{tabular}
\end{center}

\noindent The prose then cites the macro rather than the value, so tables,
figures and text share a size grid and an aggregation rule by construction
and not by inspection.

\paragraph{What Table~\ref{tab:pointwise} covers.} Five qualifications, kept
here and not in the caption, which had grown into a methods subsection
set in small type. The estimand is the schedule-pooled relative error of
Eq.~\eqref{eq:em-pooling}: per-level squared error and reference score energy
are pooled before the square root, per seed, with the per-level reference
energies of the frozen test bundle recovered deterministically by
\codefile{tools/make\_ref\_energy.py} (the bundle is shared across seeds,
so those energies are twelve numbers, validated against exp\_31's stored
sums). Both arms select their optimisation budget on a disjoint
validation bundle, the network its parameterisation \emph{and} training
length, EM--BP its iteration count, and that choice agrees with the
test-set optimum in $\netagree\%$ and $\emagree\%$ of cells respectively, so
the protocol is symmetric and the selection is close to an oracle without
being one. The $\nseq=\extendedn$ row uses a raised budget, the caps having
been set for $\nseq\le\nseqcalib$; rerunning $\nseq=\nseqcalib$ at that budget
on the same seeds moves its pooled ratio by $\budgetdelta\pm\budgetdeltase$
--- small against the ratio's own size, and in the direction that makes the
capped table conservative for the structured arm, so the caps do not carry
the comparison. The resolution gate excludes $\ndropped$ of
$\ncellstotal$ cells, because in those the narrowest fitted mixture component is
under two grid cells wide; including them changes no ratio by more than
$\dropshift$. The network wins in $\nreversed$ of the $\ncellsused$ reported
cells, both at the smallest sizes --- Appendix~\ref{app:aggregation} resolves
where the advantage is largest.

\paragraph{What Table~\ref{tab:structured} covers.} The confirmatory run scores
three regions per cell --- \structregionlist{} --- of which the table displays
two; the third is the predeclared interior slice, and it is the one the
architecture screen selects on. That is why the $\structcells$ scored cells are
$\structseeds \times \structsizes \times \structregions$ and not the
$\structseeds \times \structsizes \times 2$ a two-column table suggests. The
count is a descriptive diagnostic and not a replication count: regions are a
scoring choice applied after fitting, and sizes share nothing but the seed, so
the inferential unit behind each displayed row is the $\structseeds$ paired
seeds. Both arms are checkpointed and selected on a validation bundle disjoint
from the test bundle, and neither bundle is shared with training.

\paragraph{The architecture screen.} Table~\ref{tab:screening} reports it in
full. Its source is \texttt{screening.csv} under
\texttt{outputs/frozen/exp\_31\_screen/} at commit \texttt{569a67c}, which is the
artifact the confirmatory run names in its own recorded configuration
(\texttt{winners=auto:} that path) rather than a reconstruction. $\screenconfigs$ configurations over $\screendevseeds$
development seeds, at two training sizes, with learning rates
\screenlrs{}, both output parameterisations, and a checkpoint ladder from
$\screenckptlo$ to $\screenckpthi$ steps shared across all three
architectures. Selection pools squared error and reference energy over the
development seeds and sizes before the square root, scored on
the \screenregion{} region --- the same accumulation as the selector inside
the confirmatory job, which is what makes the screen table document the
screen that ran; the development seeds are numbered from $10000$ and
share no seed with the confirmatory set, so the reported interval is not a
selection artefact. The winner is the window head at radius
$\screenwinnerradius$, width $\screenwinnerwidth$, learning rate
$\screenwinnerlr$, \screenwinnermode{} parameterisation. Parameter counts
across the grids are \screenwindowparams{} (window), \screenbimpparams{}
(message passing) and \screenconvparams{} (dilated convolution): the arm with
the most capacity placed last, which is the relevant check on whether the
comparison was tuned evenly.

\paragraph{The convergence sweep.}
The settling times of Section~\ref{sec:em-stopping-times} come from
\texttt{outputs/frozen/}\allowbreak\texttt{exp\_27\_*}, run on 16 August
2026 under SLURM job
\texttt{628601} and its shards --- a recorded-configuration-only run in the
sense above, so the gate refuses it by default and the generator passes
\texttt{allow\_legacy}; the route disclosure above goes with it.
The quoted numbers come
from the $\convcap$-iteration shard, the only run length in the sweep at which
no configuration settles at its own cap; the $400$-iteration shard is censored
by its cap and its settling times are lower bounds, so it is used only for the
count of configurations that miss a $\convbudget$-iteration budget, which
censoring cannot inflate. The measurement replaces a figure of ``on the order of
$2000$ iterations'' that appeared in earlier drafts and was never measured: it
had been carried over from a diagnostic that ran a fit \emph{out to} $2000$
iterations to show its per-edge gain had plateaued.

\paragraph{The transition array's normalisation.} Stated in
Section~\ref{sec:em-numerical} and repeated here because it determines which
probabilistic object the M-step optimises: columns are \emph{not} renormalised
under the quadrature weights, so the fitted object is the continuum kernel read
through a quadrature rule and the surrogate carries no $\theta$-dependent
log-normaliser. The column-mass residual is the measured distance to a
normalised finite-state kernel.

\paragraph{Excluded inputs.} $\nseq=8192$ is excluded from both
Table~\ref{tab:pointwise} and Figure~\ref{fig:em-value-of-structure}: the
intersection of its dirty list with \texttt{exp\_07}'s import closure
contains the experiment, the frozen configuration and \texttt{src/em.py}
themselves, so the run is not recoverable. The $\nseq=4096$ rows are
retained: their intersection is empty, so no edit in that tree could have
moved those numbers.

\section{Environment}

Python $\geq 3.10$ with NumPy, SciPy, pandas and Matplotlib. The neural
baselines of Chapter~\ref{ch:em-results} are implemented in pure NumPy
(\texttt{src/nnet.py} for the fully connected denoiser,
\texttt{src/seq\_nets.py} and \texttt{src/local\_head.py} for the
locality-respecting architectures) with an explicit
\texttt{numpy.random.Generator} driving initialisation and minibatch order,
so training is stochastic but reproducible from the recorded seed. Two
optional dependencies are used and neither is on the exactness path: CuPy,
selected by the \texttt{BP\_DEVICE} environment variable, runs the same
single BP implementation on a GPU, \texttt{tests/test\_backend\_parity.py}
asserts the two devices agree, so a device change fails a test rather than
moving a number quietly, and PyTorch, imported lazily in
\texttt{src/fid.py} alone, for the generative metric. Arithmetic is float64
throughout.

The sweeps of Chapters~\ref{ch:em-method} and~\ref{ch:em-results} run under
SLURM from the batch scripts in \texttt{hpc/}, one per experiment. The
document compiles with a single run of \texttt{tectonic main.tex}.

Those are version floors rather than a lock. They do not pin the BLAS and
LAPACK implementations behind NumPy, the SciPy special-function versions, or
the \TeX{} bundle, and each of those can move a last digit; the parity test
above catches a device change but not a library one. Byte-level reproduction
of the frozen outputs is therefore not claimed --- only that the pipeline is
deterministic given its seeds on a fixed installation, which is what the
digests below record.

\section{Build manifest}

The consistency checks this document runs are internal to it, so they cannot
tell a reader which sources produced it. This manifest can. It is written at
build time by the generator
\codefile{tools/\allowbreak make\_build\_manifest.py}, and records the state
the artefact was made from:

\begin{center}
\begin{minipage}{\linewidth}\small
\begin{tabular}{@{}ll@{}}
\toprule
repository revision & \bmreposha{} \\
thesis subtree & last touched \bmthesissha{}, \bmthesisstate{} \\
research subtree & last touched \bmresearchsha{}, \bmresearchstate{} \\
generated inputs & \bmgenerated{} over \bmgeneratedcount{} files \\
figures & \bmfigures{} over \bmfigurecount{} files \\
previous \texttt{main.pdf} & \bmpdf{} \\
engine & \bmtectonic{} \\
\bottomrule
\end{tabular}
\end{minipage}
\end{center}

\noindent
This build was made from \bmlayout{}. The documents and the research package
live in one repository either way, so the revision is a single fact; the rows
below it record the last commit that touched each subtree and whether that
subtree carries uncommitted edits, both path-limited, so unrelated work
elsewhere in the tree flags neither. Where the repository root \emph{is} the
thesis, and in any repository whose history is a single commit, the three
values coincide --- not a degenerate manifest but an accurate one: there is
only one revision to name.

Three honest caveats. The PDF digest is necessarily that of the build before
this one, since writing the manifest changes the input to the current build;
it fixes the artefact one step back rather than pretending to hash itself. The
recorded revision is likewise one step behind in a repository with history,
since a manifest committed in a commit cannot name that commit; in a
single-commit publication it names that commit itself. And
a manifest fixes the \emph{source} side only: it says which code and which
generated inputs produced this document, not that the frozen experiment
outputs are themselves certified --- that is the route-by-route disclosure of
Section~\ref{sec:repro-provenance}, empty revision fields included.

% Undo the appendix's loosened line breaking, which would otherwise leak
% into the appendices that follow. This was \fussy, the partner of a
% \sloppy that no longer exists.
\setlength{\emergencystretch}{0pt}
